\documentclass[10pt,a4paper]{amsart}
\usepackage[margin=19mm]{geometry}
\usepackage{amsmath,amssymb,mathtools}
\usepackage[scr=rsfs,cal=euler]{mathalfa}
\usepackage{xfrac}
\usepackage[unicode,hidelinks,bookmarks=false]{hyperref}
\newcommand{\ie}{i.e.\ }
\newcommand{\cf}{cf.\ }
\newcommand{\resp}{respectively\ }
\newcommand{\Subsection}{\subsection}
\providecommand{\nobreakdash}{\nobreak\hbox{-}}
\setlength{\emergencystretch}{2em}
\multlinegap0pt
\usepackage{amssymb}
\usepackage{bm}
\usepackage{float}
\newtheorem{proposition}{Proposition}
\newcommand{\fq}{\mathbb{F}_q}
\newcommand{\mc}[1]{\mathcal{#1}}
\newcommand{\ssm}[1]{\sum_{#1}}
\newcommand{\td}[1]{\bm{#1}}
\newcommand{\tx}[1]{\text{#1}}
\title{On the Relations Among Algebraic Models for the MinRank Problem}
\author{Peigen Li, Siyong Tao}
\date{Source: arXiv:2609.07223v1 (CC BY 4.0)}
\begin{document}
\maketitle

\noindent\emph{Source and licence.} On the Relations Among Algebraic Models for the MinRank Problem. Version arXiv:2609.07223v1 (CC BY 4.0), distributed by arXiv under the Creative Commons Attribution 4.0 International licence, http://creativecommons.org/licenses/by/4.0/, as recorded in the arXiv licence metadata for this exact version and shown on the abstract page https://arxiv.org/abs/2609.07223v1. Full licence notice: \texttt{licenses/arxiv-2609.07223-CC-BY-4.0.txt}.\par\smallskip

\noindent\textit{Editorial scope.} Counted unit: the article's proposition Vmin (source0.tex lines 521--523). The article's complete original proof of it is delivered with it (source0.tex lines 525--548, one \texttt{proof} environment of 24 lines). No mathematics is written, repaired or re-derived: the statement and the proof are the article's own lines, in the article's own notation, and no retained line is substituted or re-flowed.  No further source-local context is needed: every cross-reference of every retained line resolves inside the two delivered spans.  Nothing was omitted from any retained span, so the package carries no omission record.  No retained line contains a citation, so the wrapper carries no bibliography and no bibliography entry.  No retained line contains a figure environment, an image-inclusion command or drawing code, so no image is delivered and the package declares no asset dependency.  Editorial formatting only, in addition: an AMS \texttt{amsart} preamble that restates the article's own declarations and macros used by the retained lines, namely the environments \texttt{proposition} on separate counters, together with the standard packages those lines need.  The retained environments print in this document as Proposition 1 in order of appearance, whereas the article prints the same items with the numbering of its own full document.  The complete native source of the article as distributed in the arXiv e-print for this exact version is bundled unchanged as \texttt{sources/209611/source0.tex}.
	\begin{proposition}\label{pro:Vsup and Vmin}
		$\pi(V^{\tx{Gb}}_{\tx{supp}})=V_{\tx{minors}}$.
	\end{proposition}

	\begin{proof}
		if $\td{a}\in V_{\tx{minors}}$, then $\tx{rank}\,\td{M_a}=r_0\leq r$. Choose linearly independent row vectors $\{\td{w}_i\}_{i=1}^{r_0}\in\fq^n$ such that $\{\td{w}_i\}_{i=1}^{r_0}$ span the row space of $\td{M_a}$. Extend $\{\td{w}_i\}_{i=1}^{r_0}$ to a basis $\{\td{w}_i\}_{i=1}^n$ of $\fq^n$. We write $\{\td{w}_i\}_{i=1}^r$ as the rows of the $r\times n$ matrix:
		\begin{equation*}
			W=\begin{pmatrix}
				\td{w}_1\\
				\vdots\\
				\td{w}_r
			\end{pmatrix}.
		\end{equation*}
		For each $i$, we append the $i$-th row of $\td{M_a}$ to get the $(r+1)\times n$ matrix of rank $r$: 
		\begin{equation*}
			W_i=\begin{pmatrix}
				(\td{M_a})_{i,*}\\
				W
			\end{pmatrix}.
		\end{equation*}
		Set $c_T=|W_{*,T}|$ for $T\in\mc{P}_r(n)$. $\tx{rank}\,W=r$, hence $\td{c}=(c_T)_{T\in\mc{P}_r(n)}\neq\td{0}_N$. $\ssm{j\in S}(-1)^{\tx{sgn}(j,S)}u_{i,j}(\td{a})c_{S\setminus\{j\}}$ is the $(r+1)$-minor $|(W_i)_{*,S}|$. But every $W_i$ has rank $r$, so $(\td{a},\td{c})\in V^{\tx{Gb}}_{\tx{supp}}$. $V_{\tx{minors}}\subseteq\pi(V^{\tx{Gb}}_{\tx{supp}})$.
		
		Conversely, fix $(\td{a},\td{c})\in V^{\tx{Gb}}_{\tx{supp}}$. If $c_{T_1}\neq 0$ for some $T_1\in\mc{P}_r(n)$, then columns of $\td{M_a}$ indexed by $T_1$ span the column space of $\td{M_a}$. Indeed, for any $j\notin T$, $h_{i,T_1\cup\{j\}}(\td{a},\td{c})=0$ ($i\in[m]$) implies that
		\begin{equation*}
			(\td{M_a})_{*,j}=(-1)^{\tx{sgn}(j,T_1)}c_{T_1}^{-1}\ssm{l\in T_1}(-1)^{\tx{sgn}(l,T_1\cup\{j\})}c_{T_1\cup\{j\}\setminus\{l\}}(\td{M_a})_{*,l},
		\end{equation*}
		here $(\td{M_a})_{*,j}$ is the $j$-th column of $\td{M_a}$. Thus $\tx{rank}\,\td{M_a}\leq r$. $\pi(V^{\tx{Gb}}_{\tx{supp}})=V_{\tx{minors}}$. 
	\end{proof}

\end{document}
